\section{Introduction}
\label{sec:intro}

AI agents increasingly talk to each other, and humans increasingly oversee them by reading what they say. This oversight rests on an assumption: if two agents coordinate on something harmful, the coordination will show up in the transcript. Steganography breaks this assumption. Agents can hide a message inside text that looks ordinary, and a monitor that reads the text sees nothing unusual~\citep{motwani2024secret,mathew2024hidden,zolkowski2025early}.

Classical steganography solves this problem for a sender and receiver who \emph{share a secret key}~\citep{cachin1998information,hopper2002provably}. For oversight, this would be good news. If a secret language needs a pre-shared key, then controlling how agents are set up (who provisions them, what memory they share, how they are fine-tuned) is enough to prevent it. Our main question is whether this is true. {\bf Can two agents communicate covertly if they never shared a secret key? If they can, what replaces the key?}

The literature gives two partial answers that point in different directions. On one hand, cryptography shows that keyless covert communication is possible. Public-key steganography~\citep{vonahn2004public}, the recent keyless key-exchange protocol of \citet{vaikuntanathan2026undetectable}, and the Codetta implementation~\citep{pang2026codetta} all let two parties agree on a key inside ordinary-looking LLM outputs. On the other hand, LLM evaluations find that agents without a pre-agreed scheme coordinate poorly. Two isolated Claude instances agree on a collusion signal only 3.4\% of the time~\citep{jarviniemi2025subversion}, and frontier models rarely pass a hidden message past a monitor without prior coordination~\citep{zolkowski2025early}. These views are compatible, because every cryptographic result \emph{assumes} that the agents already know the protocol. No prior work separates the key from this protocol knowledge and the other ingredients that ``shared secret'' quietly bundles together. No prior work states what secrecy a tacitly agreed code can have against a monitor that knows as much as the agents.

We fill this gap with a mostly theoretical analysis and two light experiments. We decompose ``shared secret'' into six ingredients: a key $K$, common knowledge of the protocol $P$, public randomness $\mathsf{R}$, channel entropy $H$, shared weights $W\!t$, and each agent's fresh private coins $\rho$. We then ask which subsets are enough for which goal against which monitor. The central tool is a simple simulation lemma (\secref{sec:theory}): whatever a receiver can compute from information the monitor also has, the monitor can compute with the same success probability. One consequence is a ``Kerckhoffs--Schelling'' corollary. A convention that agents select without communication, by reasoning about what is most natural or salient, can be guessed and decoded by a same-knowledge monitor exactly as well as by the partner. Focal points are focal \emph{because} they follow from shared knowledge, and the monitor shares that knowledge.

We test the two load-bearing predictions. In \eone, a \qwenbig sender hides a phrase in a helpful answer, and we compare a partner copy to a same-model auditor. Without an agreed scheme, the partner recovers the phrase 48\% of the time and the auditor 50\% (exact McNemar $p=0.75$). With a shared scheme, the partner beats an unaware auditor (21\% vs.\ 1\%), but an auditor told the scheme does at least as well as the partner (32\%). In \etwo, we implement a toy keyless covert key exchange through the sampler of \qwenchan. Two agents agree on a 254-bit Diffie--Hellman key in 10/10 sessions and deliver a payload. An auditor that knows the model, prompts, protocol, and public randomness decodes 0/10 payloads and cannot tell the transcripts from honest ones (NLL AUROC 0.46). We also find a new failure mode. When the channel embeds bits only at high-entropy positions, a warden that resamples just 0.25\% of tokens drops key agreement from 10/10 to 2/10. The warden does not flip many bits. It inserts and deletes them, so the bit stream falls out of sync.

Our contributions are:
\begin{itemize}[leftmargin=*,itemsep=0pt,topsep=0pt]
    \item We formalize ``shared secret'' as six separable ingredients and prove which ones are necessary for coordination, undetectability, and secrecy against weaker, equally informed, and active monitors (\tabref{tab:necessity}).
    \item We prove that tacit (Schelling) conventions and one-way keyless schemes give exactly zero advantage over an equally informed monitor, and that an undetectable protocol cannot be learned from the channel. The real shared secret is prior common knowledge of the protocol, which may be public.
    \item We conduct an end-to-end test with open-weight LLMs showing no colluder advantage over a same-model auditor when no scheme is agreed, and no advantage once the auditor knows a shared scheme.
    \item We implement a working keyless covert key exchange through a real LM sampler, measure its cost (about 7 tokens per key bit), and identify desynchronization under active resampling as a cheap defense against entropy-gated channels.
\end{itemize}

The rest of the paper is organized as follows. \Secref{sec:related} reviews related work. \Secref{sec:theory} presents the model and our theoretical results. \Secref{sec:setup} describes the experiments, and \secref{sec:results} reports them. \Secref{sec:discussion} discusses implications and limitations.
